% ==================== 致谢 ====================
{
\ctexset{chapter/format = \centering\zihao{3}\heiti\bfseries}
\chapter*{致\hspace{2em}谢}
\addcontentsline{toc}{chapter}{致谢}

\zihao{-4}\songti
本论文从选题到定稿历时近一年，最先想感谢的是我的导师孙成娇老师。开题时我在若干方向之间摇摆不定，孙老师帮我把题目收敛到Apollo规划模块的改进；初稿有不少含糊其辞的地方，孙老师逐段批注，也多次提醒我不要急着下结论、先把实验跑扎实。这些意见让我开始明白学术写作与写代码是两回事。计算机工程学院几位任课老师四年来在算法与工程课上的训练，是我能够读进Apollo源码的底子，一并致以谢意。

百度Apollo团队把一套工业级自动驾驶栈完整开源，我才有机会把本科阶段的改进想法真正落到一个可运行的系统上；Apollo开发者社区里围绕PathBoundsDecider与SpeedBoundsDecider的若干讨论帖，多次帮我走出卡壳。

同门与室友陪我熬过若干次通宵调参与推公式的夜晚，彼此吐槽也彼此鼓劲。最后要写给我的家人，这一年我很少回家，他们从不催问进度，只说注意身体，这份沉默的托底是我真正的底气。
}
